\documentclass[aspectratio=169, 14pt]{beamer}
\usepackage{beamer_preamble}

\begin{document}

% ==== Title ===============================================================
\titleslide{Lesson 5-4 \textperiodcentered\ Period 3}{Radical Equations --- Assessment-Critical: Isolate a Variable + Table Completion}

% ==== Do Now ==============================================================
\begin{frame}{Do Now --- Extraneous Solutions (HOT)}{Algebra 2 \textperiodcentered\ Lesson 5-4 \textperiodcentered\ Period 3}
\phasebadges{3}{5}{bridge from P2}
\vspace{0.2cm}
\begin{projcard}{Higher Order Thinking \textperiodcentered\ Extraneous Solutions}{slidegold}
\small
Some, but not all, equations with rational exponents have extraneous solutions.

\vspace{0.2cm}
\textbf{Think:} What is the relationship between the exponents and the possibility
of having extraneous solutions? Explain your reasoning.

\vspace{0.2cm}
\textcolor{slidegold}{\textbf{Predict:}} If you solved \(V = \tfrac{1}{3}\pi r^2 h\) for \(r\),
would you expect an extraneous solution? Why or why not?
\end{projcard}
\end{frame}

% ==== Launch ==============================================================
\begin{frame}{Launch --- Rewrite a Formula + Table Completion}{Algebra 2 \textperiodcentered\ Lesson 5-4 \textperiodcentered\ Period 3}
\phasebadges{2}{12}{teacher models Ex~2 + \#44}
\vspace{0.15cm}
\begin{projcard}{Isolate-a-Variable Rules}{slidegreen}
\small
\begin{enumerate}
  \setlength{\itemsep}{1pt}
  \item Identify which variable to isolate.
  \item Undo operations LAST-IN, FIRST-OUT: add/sub \(\rightarrow\) mult/div \(\rightarrow\) exponent/radical.
  \item Raise both sides to the \textbf{RECIPROCAL} power to undo \(x^{m/n}\): flip to \(x^{n/m}\).
\end{enumerate}
\end{projcard}
\vspace{-0.1cm}
\begin{projcard}{Table-Evaluation Rules}{slidegreen}
\small
\begin{enumerate}
  \setlength{\itemsep}{1pt}
  \item Substitute each input into the rewritten formula.
  \item Simplify rational exponents step-by-step (use calculator for irrational outputs).
  \item Round to specified precision.
\end{enumerate}
\end{projcard}
\end{frame}

% ==== Launch card 2: Example 2 + #44 =====================================
\begin{frame}{Launch --- Ex 2 (Golden Gate) + Practice \#44}{Algebra 2 \textperiodcentered\ Lesson 5-4 \textperiodcentered\ Period 3}
\phasebadges{2}{12}{teacher models}
\vspace{0.15cm}
\begin{projcard}{Example 2 \textperiodcentered\ Rewrite the Golden Gate formula}{slidegreen}
\small
Given: \(x = \dfrac{\sqrt{y - 220}}{0.010583}\) \quad Isolate \(y\).

\vspace{0.1cm}
Steps: multiply both sides by coefficient \(\rightarrow\) square both sides
\(\rightarrow\) add 220.\\
Then substitute \(x = 200\) and evaluate \(y\).
\end{projcard}
\vspace{-0.1cm}
\begin{projcard}{Practice \#44 \textperiodcentered\ Table Completion}{slideaccent}
\small
\textcolor{slideaccent}{\bfseries ASSESSMENT REHEARSAL --- Item 2}

Complete the table for \(y = \sqrt[3]{2x + z} - 12\).
Teacher walks through rows 1--2 live. Copy the steps; note calculator workflow for irrational outputs.
\end{projcard}
\end{frame}

% ==== Explore =============================================================
\begin{frame}{Explore --- Think-Pair-Share (35 min)}{Algebra 2 \textperiodcentered\ Lesson 5-4 \textperiodcentered\ Period 3}
\phasebadges{2}{35}{student work --- TPS}
\small\textcolor{slidegray}{6 items in order. Never cut \#10 (Item 1A rehearsal). Cut \#27 first if pace slows.}

\vspace{0.1cm}
\begin{projcard}{Try It 2 \textperiodcentered\ Stopping distance formula}{slideblue}
\(v = 3.57\sqrt{d}\) \quad Rewrite to isolate \(d\).
\end{projcard}
\vspace{-0.12cm}
\begin{projcard}{Practice \#10 \textperiodcentered\ Isolate \(x\)}{slideaccent}
\textcolor{slideaccent}{\bfseries ASSESSMENT ITEM 1A REHEARSAL} \quad
Rewrite \(y = \sqrt{(x-48)/6}\) to isolate \(x\).
\end{projcard}
\vspace{-0.12cm}
\begin{projcard}{Practice \#25 / \#26 / \#27 \textperiodcentered\ Rationalization pattern}{slideblue}
\(\#25\!:\ x=3\sqrt[3]{15+y}\) \quad
\(\#26\!:\ x=\tfrac{\sqrt{2y}}{26}\) \quad
\(\#27\!:\ x=\tfrac{\sqrt{y-14.2}}{0.05}\) \quad Solve for \(y\).
\end{projcard}
\vspace{-0.12cm}
\begin{projcard}{Practice \#39 \textperiodcentered\ BSA Application (modeling)}{slideblue}
\(BSA = \sqrt{\tfrac{H \cdot M}{3600}}\), \(BSA < 2.0\), \(H = 165\) cm.
Find max mass (kg) to nearest tenth.
\end{projcard}
\end{frame}

% ==== Share / Summary =====================================================
\begin{frame}{Share / Summary}{Algebra 2 \textperiodcentered\ Lesson 5-4 \textperiodcentered\ Period 3}
\phasebadges{2}{5}{concept summary + assessment preview}
\vspace{0.15cm}
\begin{projcard}{Sentence Frame (\#10 or \#39)}{slideblue}
\small
``To isolate \underline{\hspace{0.7cm}}, I first \underline{\hspace{1cm}}, then raise both sides to the power
\underline{\hspace{0.7cm}}. The rewritten formula is \underline{\hspace{1cm}}.
Plugging in \underline{\hspace{0.5cm}} = \underline{\hspace{0.5cm}} gives \underline{\hspace{0.7cm}}.
In context this means \underline{\hspace{1.2cm}}.''
\end{projcard}
\vspace{-0.1cm}
\begin{projcard}{Topic 5 Assessment Preview}{slideaccent}
\small
\textcolor{slideaccent}{\bfseries Items 1A/1B} --- Rewrite a pyramid-volume formula for \(x\), then evaluate.
\(\Rightarrow\) Rehearsed by \textbf{Practice \#10}.

\textcolor{slideaccent}{\bfseries Item 2} --- Complete the table from the rewritten formula.
\(\Rightarrow\) Rehearsed by \textbf{Practice \#44} (done in Launch).
\end{projcard}
\end{frame}

% ==== Exit Ticket =========================================================
\begin{frame}{Exit Ticket --- Pre-Assessment Confidence Check}{Algebra 2 \textperiodcentered\ Lesson 5-4 \textperiodcentered\ Period 3}
\phasebadges{1-2}{3}{not graded}
\vspace{0.2cm}
\begin{projcard}{Complete on your exit ticket}{slidegold}
\begin{enumerate}
  \setlength{\itemsep}{0.25cm}
  \item To isolate a variable from a rational-exponent formula I \underline{\hspace{4cm}}.
  \item When I evaluate the formula at a specific \(x\), I \underline{\hspace{4cm}}.
  \item One biggest thing I learned today: \underline{\hspace{4.5cm}}.
\end{enumerate}
\end{projcard}
\end{frame}

\end{document}
